%% toc: P1 1 1.5
\section{Proofs for changing the formal language}
\label{sec:annex-p01-w02-p01-003-language-change}

Throughout, \(\Lang'\) extends \(\Lang\) unless a statement says otherwise,
and the notation is that of
\cref{def:p01-w02-p01-003-signature-subscripts}.

\begin{annexitem}{annex:p01-w02-p01-003-syntax-inclusion}{thm:p01-w02-p01-003-syntax-inclusion}{The old syntax inside the new}
\begin{proof}
\emph{Clause 1.} \(\operatorname{Term}(\Lang)\) is the least set that
contains every variable and contains \(f(t_1,\ldots,t_n)\) whenever
\(f\in\mathcal F\) has arity \(n\) and \(t_1,\ldots,t_n\) belong to it. The
set \(\operatorname{Term}(\Lang')\) contains every variable. If
\(f\in\mathcal F\) has \(\Lang\)-arity \(n\), then \(f\in\mathcal F'\) and its
\(\Lang'\)-arity is also \(n\), because the arity map of \(\Lang'\) restricts
to that of \(\Lang\); so \(\operatorname{Term}(\Lang')\) contains
\(f(t_1,\ldots,t_n)\) whenever it contains \(t_1,\ldots,t_n\). Thus
\(\operatorname{Term}(\Lang')\) satisfies both closure conditions defining
\(\operatorname{Term}(\Lang)\), and leastness gives
\(\operatorname{Term}(\Lang)\subseteq\operatorname{Term}(\Lang')\).

The same argument applies to formulas. \(\operatorname{Fm}(\Lang)\) is the
least set containing \(r=t\) for \(r,t\in\operatorname{Term}(\Lang)\),
\(R(t_1,\ldots,t_n)\) for \(R\in\mathcal R\) of arity \(n\) and
\(t_i\in\operatorname{Term}(\Lang)\), and \(\bottom\), and closed under
\(\neg\), the binary connectives \(\land,\lor,\to\) and both quantifiers
over every variable. By the term inclusion just proved, every such identity
is in \(\operatorname{Fm}(\Lang')\). Every \(R\in\mathcal R\) lies in
\(\mathcal R'\) with the same arity, so every such relational atom is in
\(\operatorname{Fm}(\Lang')\). Also \(\bottom\in\operatorname{Fm}(\Lang')\),
and \(\operatorname{Fm}(\Lang')\) has the same closure under connectives and
quantifiers. Leastness gives
\(\operatorname{Fm}(\Lang)\subseteq\operatorname{Fm}(\Lang')\).

\emph{Clause 2.} The recursive clauses defining \(\operatorname{Var}\),
\(\operatorname{FV}\), substitution in a term and capture-avoiding
substitution in a formula mention no signature symbol except as material
copied from the input. The one choice made in them is the fresh variable in
the renaming clause, which is the least variable outside
\(\operatorname{Var}(\psi)\cup\operatorname{Var}(t)\cup\{x,y\}\) in the fixed
order of \(\mathsf{Var}\) \cite[ch.~``First-Order Syntax'', section
``Capture-avoiding substitution'']{FrozenFormalLogicLanguageChange}. Hence for
inputs drawn from \(\Lang\) the two computations, over \(\Lang\) and over
\(\Lang'\), apply the same clause at each step to the same strings.

To make this precise, first take terms. For \(r,t\in\operatorname{Term}(\Lang)\)
we show by induction on \(r\) that \(r[t/x]\) is the same string over both
signatures and lies in \(\operatorname{Term}(\Lang)\). If \(r\) is \(x\), the
result is \(t\); if \(r\) is a variable \(y\ne x\), it is \(y\). If
\(r=f(r_1,\ldots,r_n)\) with \(f\in\mathcal F\), the result is
\(f(r_1[t/x],\ldots,r_n[t/x])\); by the induction hypothesis the arguments
are the same strings over both signatures and lie in
\(\operatorname{Term}(\Lang)\), so the result does too. The same induction
shows that \(\operatorname{Var}(r)\) is the same set over both signatures,
since its clauses are a singleton for a variable and a union over the
arguments.

Next take \(\varphi\in\operatorname{Fm}(\Lang)\) and proceed by strong
induction on the complexity \(c(\varphi)\), simultaneously for all
\(t\in\operatorname{Term}(\Lang)\) and \(x\in\mathsf{Var}\). For an atomic
\(\varphi\) the substitution acts on its terms, and the term case applies;
the result is an atom over the same symbol, hence in
\(\operatorname{Fm}(\Lang)\). For \(\bottom\) the result is \(\bottom\). The
connective clauses commute with substitution and use the induction
hypothesis on the components. For \(\varphi=Qy\,\psi\):
\begin{itemize}
  \item if \(y=x\), the result is \(\varphi\) itself;
  \item if \(y\ne x\) and \(y\notin\operatorname{Var}(t)\), the result is
        \(Qy\,(\psi[t/x])\), and the induction hypothesis applies to
        \(\psi\);
  \item if \(y\ne x\) and \(y\in\operatorname{Var}(t)\), the variable \(z\)
        is determined by \(\operatorname{Var}(\psi)\),
        \(\operatorname{Var}(t)\), \(x\) and \(y\), which are the same over
        both signatures by the induction hypothesis and the term case. The
        induction hypothesis for \(\psi\) and the variable term \(z\) gives
        that \(\psi[z/y]\) is the same string over both signatures and lies
        in \(\operatorname{Fm}(\Lang)\). Its complexity equals \(c(\psi)\)
        by \cref{lem:substitution-complexity}, which is below
        \(c(\varphi)\), so the induction hypothesis applies to
        \((\psi[z/y])[t/x]\) as well.
\end{itemize}
In each case the result is the same string over both signatures and lies in
\(\operatorname{Fm}(\Lang)\). The equality of \(\operatorname{Var}(\varphi)\)
and of \(\operatorname{FV}(\varphi)\) over both signatures follows by the same
induction, because their clauses for atoms use only the variables of the
terms, their clauses for connectives are unions, and their quantifier
clauses add or remove the bound variable.

\emph{Clause 3.} Write \(\equiv_\alpha\) for alpha-equivalence on
\(\operatorname{Fm}(\Lang)\) and \(\equiv'_\alpha\) for alpha-equivalence on
\(\operatorname{Fm}(\Lang')\). Each is the least equivalence relation on its
formula set that is a congruence for every formula constructor and contains
every generating renaming \(Qx\,\chi\) versus \(Qy\,(\chi[y/x])\) with
\(y\notin\operatorname{Var}(\chi)\)
\cite[ch.~``First-Order Syntax'', definition
``Alpha-equivalence'']{FrozenFormalLogicLanguageChange}. Put
\[
 E\defeq\{(\varphi,\psi)\in\operatorname{Fm}(\Lang)\times\operatorname{Fm}(\Lang):
          \varphi\equiv'_\alpha\psi\}.
\]
\(E\) is reflexive, symmetric and transitive on \(\operatorname{Fm}(\Lang)\),
because \(\equiv'_\alpha\) has these properties on the larger set. It is a
congruence for the constructors of \(\operatorname{Fm}(\Lang)\): if
\(\varphi\mathrel E\psi\) then \(\neg\varphi\equiv'_\alpha\neg\psi\) by the
congruence property of \(\equiv'_\alpha\), and \(\neg\varphi\),
\(\neg\psi\) are \(\Lang\)-formulas, so \(\neg\varphi\mathrel E\neg\psi\);
the binary connectives and the two quantifiers are handled in the same way,
and the atomic constructors take no formula arguments. Finally, let
\(\chi\in\operatorname{Fm}(\Lang)\) and \(y\notin\operatorname{Var}(\chi)\).
By clause 2, \(\chi[y/x]\) computed over \(\Lang\) equals the string computed
over \(\Lang'\) and lies in \(\operatorname{Fm}(\Lang)\), and
\(\operatorname{Var}(\chi)\) is the same over both signatures. So the pair
\((Qx\,\chi,\;Qy\,(\chi[y/x]))\) is a generating renaming of
\(\equiv'_\alpha\) whose members are both \(\Lang\)-formulas, and it belongs
to \(E\). Hence \(E\) is an equivalence relation and congruence on
\(\operatorname{Fm}(\Lang)\) containing every generating renaming of
\(\equiv_\alpha\), and leastness gives \({\equiv_\alpha}\subseteq E\). That is,
\(\varphi\equiv_\alpha\psi\) implies \(\varphi\equiv'_\alpha\psi\).
\end{proof}
\end{annexitem}

\begin{annexitem}{annex:p01-w02-p01-003-expansion}{thm:p01-w02-p01-003-expansion}{Every structure of the old signature has an expansion}
\begin{proof}
Let \(\mathcal M\) be an \(\Lang\)-structure with domain \(M\). The domain
is nonempty, so fix one \(a\in M\); this is a single instance of
existential instantiation and uses no choice principle. Define
\(\mathcal M'\) with domain \(M\) as follows:
\begin{itemize}
  \item each symbol of \(\Lang\) receives its interpretation in
        \(\mathcal M\);
  \item each function symbol \(f\in\mathcal F'\setminus\mathcal F\) of
        arity \(n\) receives the constant function
        \(M^{n}\to M\) with value \(a\); for \(n=0\) this is the element
        \(a\);
  \item each relation symbol \(R\in\mathcal R'\setminus\mathcal R\)
        receives the empty relation \(\varnothing\subseteq M^{n}\), where
        \(n\) is its arity.
\end{itemize}
These clauses are well defined because the arity of an old symbol is the
same in both signatures, so the interpretation copied from \(\mathcal M\)
has the arity \(\Lang'\) requires. Every symbol of \(\Lang'\) receives
exactly one interpretation of the right kind, so \(\mathcal M'\) is an
\(\Lang'\)-structure. It has the domain of \(\mathcal M\) and agrees with it
on every symbol of \(\Lang\), so it is an expansion of \(\mathcal M\).

Let \(\mathcal M''\) be any \(\Lang'\)-expansion of \(\mathcal M\). Its
reduct \(\mathcal M''\mathbin{\upharpoonright}\Lang\) has the domain of
\(\mathcal M''\), which is \(M\), and gives each symbol of \(\Lang\) its
interpretation in \(\mathcal M''\), which is its interpretation in
\(\mathcal M\). An \(\Lang\)-structure is fixed by its domain and its
interpretations of the symbols of \(\Lang\), so
\(\mathcal M''\mathbin{\upharpoonright}\Lang=\mathcal M\).
\end{proof}
\end{annexitem}

\begin{annexitem}{annex:p01-w02-p01-003-semantic}{thm:p01-w02-p01-003-semantic-comparison}{Semantic consequence is unchanged on the old formulas}
\begin{proof}
Let \(\Gamma\cup\{\varphi\}\subseteq\operatorname{Fm}(\Lang)\). By
\cref{thm:p01-w02-p01-003-syntax-inclusion} these are also
\(\Lang'\)-formulas, so both sides of the equivalence are defined. Each use
of \cref{lem:reduct-invariance} below applies to an \(\Lang\)-formula,
which is the hypothesis of that lemma.

\emph{Left to right.} Assume \(\Gamma\entails_{\Lang}\varphi\). Let
\(\mathcal M'\) be an \(\Lang'\)-structure and \(s\) an assignment in its
domain with \(\mathcal M',s\entails\gamma\) for every \(\gamma\in\Gamma\).
The reduct \(\mathcal M'\mathbin{\upharpoonright}\Lang\) is an
\(\Lang\)-structure with the same domain, so \(s\) is an assignment in it.
By \cref{lem:reduct-invariance},
\(\mathcal M'\mathbin{\upharpoonright}\Lang,s\entails\gamma\) for every
\(\gamma\in\Gamma\). The assumption gives
\(\mathcal M'\mathbin{\upharpoonright}\Lang,s\entails\varphi\), and
\cref{lem:reduct-invariance} again gives \(\mathcal M',s\entails\varphi\).
Since \(\mathcal M'\) and \(s\) were arbitrary,
\(\Gamma\entails_{\Lang'}\varphi\).

\emph{Right to left.} Assume \(\Gamma\entails_{\Lang'}\varphi\). Let
\(\mathcal M\) be an \(\Lang\)-structure and \(s\) an assignment in its
domain with \(\mathcal M,s\entails\gamma\) for every \(\gamma\in\Gamma\).
By \cref{thm:p01-w02-p01-003-expansion} there is an \(\Lang'\)-expansion
\(\mathcal M'\) with \(\mathcal M'\mathbin{\upharpoonright}\Lang=\mathcal M\);
its domain is that of \(\mathcal M\), so \(s\) is an assignment in it. For
each \(\gamma\in\Gamma\), \(\mathcal M'\mathbin{\upharpoonright}\Lang,s\entails\gamma\),
so \cref{lem:reduct-invariance} gives \(\mathcal M',s\entails\gamma\). The
assumption gives \(\mathcal M',s\entails\varphi\), and
\cref{lem:reduct-invariance} gives
\(\mathcal M,s=\mathcal M'\mathbin{\upharpoonright}\Lang,s\entails\varphi\).
Hence \(\Gamma\entails_{\Lang}\varphi\).

The validity statement is the case \(\Gamma=\varnothing\).
\end{proof}
\end{annexitem}

\begin{annexitem}{annex:p01-w02-p01-003-derivability}{thm:p01-w02-p01-003-derivability-comparison}{Derivability on the old formulas}
\begin{proof}
\emph{Clause 1.} Let \(D\) be a derivation in
\(\mathsf{ND}_{\mathsf{FOL}}[\Lang]\) with conclusion \(\varphi\) and open
assumptions in \(\Gamma\). It is a finite rooted tree whose nodes carry
\(\Lang\)-formulas, whose assumption leaves may carry discharge labels, and
whose inner nodes are instances of the rule schemata. The same tree, with the
same labels, is a derivation in \(\mathsf{ND}_{\mathsf{FOL}}[\Lang']\). Every
node formula is an \(\Lang'\)-formula by clause 1 of
\cref{thm:p01-w02-p01-003-syntax-inclusion}. It remains to check that each
inner node is an instance of the same rule over \(\Lang'\). Which assumption
occurrences are open at a node, and which a rule application discharges, is
fixed by the tree and its labels and is therefore the same in both readings.
\begin{itemize}
  \item The rules of \(\mathsf{ND}_{\mathsf{CPL}}\) are instances when the
        premises, conclusion and discharged assumptions have the displayed
        shapes, which are conditions of string identity. They hold over
        \(\Lang'\) because they hold over \(\Lang\) for the same strings.
  \item \(\forall E\) and \(\exists I\): the premise or conclusion has the
        form \(\chi[t/x]\) with \(t\in\operatorname{Term}(\Lang)\). By
        clause 2 of \cref{thm:p01-w02-p01-003-syntax-inclusion}, the string
        \(\chi[t/x]\) is the same over \(\Lang'\), and
        \(t\in\operatorname{Term}(\Lang')\) by clause 1.
  \item \(\forall I\): the side condition is that \(x\) is not free in any
        open assumption. The open assumptions are the same, and their
        free-variable sets are the same by clause 2.
  \item \(\exists E\): the assumption discharged is \(\chi[y/x]\), the
        same string over \(\Lang'\) by clause 2. The side condition is that
        \(y\) is absent from \(\psi\), from every other open assumption and
        from \(\exists x\,\chi\); occurrence of a variable in a given string
        does not depend on the signature.
  \item \(\alpha\): the side condition \(\chi\equiv_\alpha\chi'\) over
        \(\Lang\) yields \(\chi\equiv'_\alpha\chi'\) over \(\Lang'\) by
        clause 3.
  \item \(=I\) and \(=E\): the terms involved lie in
        \(\operatorname{Term}(\Lang)\subseteq\operatorname{Term}(\Lang')\),
        and the substitution instances \(\chi[r/x]\) and \(\chi[t/x]\) are
        the same strings over \(\Lang'\) by clause 2.
\end{itemize}
These are all the rules of the calculus
\cite[ch.~``First-Order Deduction'', section ``Natural-deduction rules for
quantifiers and equality'']{FrozenFormalLogicLanguageChange}. Hence \(D\) is a
derivation in \(\mathsf{ND}_{\mathsf{FOL}}[\Lang']\) with conclusion
\(\varphi\) and the same open assumptions, all in \(\Gamma\), so
\(\Gamma\derives_{\mathsf{ND}_{\mathsf{FOL}}[\Lang']}\varphi\).

\emph{Clause 2.} Assume the Axiom of Choice and
\(\Gamma\derives_{\mathsf{ND}_{\mathsf{FOL}}[\Lang']}\varphi\). The set
\(\Gamma\cup\{\varphi\}\) is contained in \(\operatorname{Fm}(\Lang')\), so
\cref{thm:fol-soundness}, which holds for every signature and is applied to
\(\Lang'\), gives \(\Gamma\entails_{\Lang'}\varphi\). Because
\(\Gamma\cup\{\varphi\}\subseteq\operatorname{Fm}(\Lang)\),
\cref{thm:p01-w02-p01-003-semantic-comparison} gives
\(\Gamma\entails_{\Lang}\varphi\). \Cref{thm:fol-completeness}, applied to
\(\Lang\) under the Axiom of Choice, gives
\(\Gamma\derives_{\mathsf{ND}_{\mathsf{FOL}}[\Lang]}\varphi\).
\end{proof}
\end{annexitem}

\begin{annexitem}{annex:p01-w02-p01-003-shared}{thm:p01-w02-p01-003-common-subsignature}{Comparison through a common sub-signature}
\begin{proof}
\(\Lang_0\) is a signature: its function and relation symbols are symbols of
\(\Lang_1\), so they are pairwise disjoint from the variables, from the
logical and punctuation tokens and from each other, as the definition of a
signature requires; and each receives the single arity on which
\(\Lang_1\) and \(\Lang_2\) agree. Its function symbols form a subset of
\(\mathcal F_1\) and of \(\mathcal F_2\), its relation symbols a subset of
\(\mathcal R_1\) and of \(\mathcal R_2\), and on each of them the arity maps
of \(\Lang_1\) and of \(\Lang_2\) take the value assigned by \(\Lang_0\).
Thus both \(\Lang_1\) and \(\Lang_2\) extend \(\Lang_0\).

For \(\Gamma\cup\{\varphi\}\subseteq\operatorname{Fm}(\Lang_0)\),
\cref{thm:p01-w02-p01-003-semantic-comparison} applied to the extension
\(\Lang_0\subseteq\Lang_1\) gives
\(\Gamma\entails_{\Lang_0}\varphi\) iff \(\Gamma\entails_{\Lang_1}\varphi\),
and applied to \(\Lang_0\subseteq\Lang_2\) gives
\(\Gamma\entails_{\Lang_0}\varphi\) iff \(\Gamma\entails_{\Lang_2}\varphi\).
Composing the two equivalences gives the displayed chain. Under the Axiom of
Choice, both clauses of
\cref{thm:p01-w02-p01-003-derivability-comparison} apply to each extension, so
\(\Gamma\derives_{\mathsf{ND}_{\mathsf{FOL}}[\Lang_0]}\varphi\) iff
\(\Gamma\derives_{\mathsf{ND}_{\mathsf{FOL}}[\Lang_i]}\varphi\) for
\(i=1,2\), and the derivability chain follows in the same way.
\end{proof}
\end{annexitem}

\begin{annexitem}{annex:p01-w02-p01-003-restricted}{thm:p01-w02-p01-003-restricted-class}{A new constant with an axiom that is not conservative}
\begin{proof}
\emph{Clause 1.} Let \(\mathcal M'\in\mathcal K_A\), let \(s\) be an
assignment in it, and put \(a\defeq c^{\mathcal M'}\). We show
\(\mathcal M',s[y\mapsto a]\entails\forall x\,(x=y\lor R(y,x))\). Let \(b\)
be any element of the domain and put \(s'\defeq s[y\mapsto a][x\mapsto b]\).
Since \(\mathcal M'\entails A\) and \(A\) is a sentence,
\(\mathcal M',s[x\mapsto b]\entails x=c\lor R(c,x)\), so either \(b=a\) or
\((a,b)\in R^{\mathcal M'}\). Because \(x\ne y\), the assignment \(s'\) sends
\(y\) to \(a\) and \(x\) to \(b\). In the first case
\(\mathcal M',s'\entails x=y\); in the second
\(\mathcal M',s'\entails R(y,x)\). Either way
\(\mathcal M',s'\entails x=y\lor R(y,x)\). As \(b\) was arbitrary, the
universal clause gives the claim, and the existential clause, with witness
\(a\), gives \(\mathcal M',s\entails\theta\). Hence
\(\varnothing\entails_{\mathcal K_A}\theta\).

\emph{Clause 2.} Let \(\mathcal N\) have domain \(\{0,1\}\) and
\(R^{\mathcal N}=\varnothing\), and let \(s\) be any assignment. For
\(a\in\{0,1\}\), put \(b\defeq1-a\). Under
\(s[y\mapsto a][x\mapsto b]\) the identity \(x=y\) is false because
\(b\ne a\), and \(R(y,x)\) is false because \(R^{\mathcal N}\) is empty.
So \(\mathcal N,s[y\mapsto a]\not\entails\forall x\,(x=y\lor R(y,x))\)
for both values of \(a\), and \(\mathcal N,s\not\entails\theta\). Since
\(\mathcal N\) is an \(\Lang\)-structure, \(\varnothing\not\entails_{\Lang}\theta\).

\emph{Clause 3.} Let \(\mathcal N'\) be any \(\Lang'\)-expansion of
\(\mathcal N\). Then \(R^{\mathcal N'}=\varnothing\) and
\(a\defeq c^{\mathcal N'}\in\{0,1\}\). With \(b\defeq1-a\), the formula
\(x=c\lor R(c,x)\) fails under any assignment sending \(x\) to \(b\), for the
reasons given in clause 2. So \(\mathcal N'\not\entails A\) and
\(\mathcal N'\notin\mathcal K_A\).

\emph{Clause 4.} Put \(\chi\defeq\forall x\,(x=y\lor R(y,x))\), so that
\(\theta=\exists y\,\chi\). Compute \(\chi[c/y]\) by the quantifier clause of
capture-avoiding substitution: the binder is \(x\), which differs from \(y\),
and \(x\notin\operatorname{Var}(c)=\varnothing\), so the substitution passes
under the quantifier and replaces the two free occurrences of \(y\), giving
\(\forall x\,(x=c\lor R(c,x))\), which is \(A\). The tree with the single
assumption leaf \(A\) and one application of \(\exists I\), with premise
\(\chi[c/y]\) and conclusion \(\exists y\,\chi\), is therefore a derivation
in \(\mathsf{ND}_{\mathsf{FOL}}[\Lang']\); the term \(c\) lies in
\(\operatorname{Term}(\Lang')\). Its only open assumption is \(A\), so
\(\{A\}\derives_{\mathsf{ND}_{\mathsf{FOL}}[\Lang']}\theta\).
\end{proof}
\end{annexitem}
